% TOP_PROBLEM: {"id": 3310, "title": "Chromatic number of random lifts of complete graphs", "category_id": 3, "category": {"id": 3, "name": "graph_theory", "display_name": "Graph Theory", "description": "Problems involving graphs, networks, and their properties.", "slug": "graph-theory", "order_index": 3, "created_at": "2026-07-31T15:26:25.671Z"}, "status": "open", "problem_number": "OPG-37656", "set_id": 12}
% FIRST_POSTED: 2026-10-01
% ENTRY_KIND: statement_only
% UnsolvedMath ID 3310; source catalogue code OPG-37656
% !TeX program = lualatex
% Definition and sources only; this file is not an LLM solution attempt.
\documentclass[11pt,a4paper]{article}
\usepackage[margin=25mm]{geometry}
\usepackage{fontspec}
\setmainfont{lmroman10-regular.otf}[BoldFont=lmroman10-bold.otf,
 ItalicFont=lmroman10-italic.otf,BoldItalicFont=lmroman10-bolditalic.otf]
\usepackage{amsmath,amssymb,mathtools}
\usepackage{unicode-math}
\setmathfont{latinmodern-math.otf}
\AtBeginDocument{\let\setminus\smallsetminus}
\usepackage{xurl,hyperref}
\hypersetup{colorlinks=true,urlcolor=blue}
\setcounter{secnumdepth}{0}
\setlength{\parindent}{0pt}
\setlength{\parskip}{5pt}
\setlength{\emergencystretch}{3em}
\begin{document}
\section*{Chromatic number of random lifts of complete graphs}\label{3310}
{\small UnsolvedMath ID 3310 · OPG-37656 · Graph Theory · OpenGarden}
\subsection{Definitions and mathematical statement}
For an integer $h\ge1$, a labelled $h$-lift of the complete graph $K_5$
has vertex set $\{1,2,3,4,5\}\times\{1,\ldots,h\}$. For each pair
$1\le i<j\le5$, choose a bijection
$\pi_{ij}:\{1,\ldots,h\}\to\{1,\ldots,h\}$ and put in the $h$ edges
\[
                       (i,a)(j,\pi_{ij}(a)),\qquad 1\le a\le h.
\]
There are no other edges. Each of the five fibres is independent,
and the edges between each pair of fibres form a perfect matching.
Every vertex therefore has degree four.

Let $L_h$ be the random lift obtained by choosing the ten bijections
independently and uniformly from all $h!$ permutations. This defines
the probability model on labelled lifts; it is not uniform over
unlabelled isomorphism classes. Write $\chi(L_h)$ for the minimum
number of colours in a proper vertex colouring.

The question is whether there exists a fixed integer $q$ such that
\[
                     \lim_{h\to\infty}
                     \mathbb P\bigl(\chi(L_h)=q\bigr)=1.
\]
If so, determine the value $q$. The base graph is fixed as $K_5$
throughout; only the size of the fibres tends to infinity. One asks
for exact concentration on a single integer, rather than bounded
variance, concentration in an interval of several integers, or a
limit for the expected chromatic number. A colouring using fibre
labels always uses five colours, but its existence alone gives no
answer to this single-value concentration question.
\subsection{Short English statement}
As the fibre size grows, does the chromatic number of a random lift of K5 equal one fixed integer with probability tending to one?
\subsection{Sources}
\begin{itemize}
\item[S1] ULAM.AI and the UnsolvedMath Contributors, supplied problems.json, record 3310 (OPG-37656), OpenGarden collection. Catalogue identity and source transcription; CC BY 4.0.
\url{https://huggingface.co/datasets/ulamai/UnsolvedMath}.
\item[S2] Open Problem Garden, Chromatic number of random lifts of complete graphs. Original problem and discussion; the mathematical formulation below is expanded from the supplied source record.
\url{http://www.openproblemgarden.org/op/chromatic_number_of_random_lifts_of_complete_graphs}.
\end{itemize}
\end{document}
